\begin{description}
    \item[\textbf{Acurácia}] Métrica que expressa a proporção de sinais corretamente classificados pelo modelo em relação ao total de amostras testadas, representada em percentual.

    \item[\textbf{Aprendizado Profundo}] Subcampo do aprendizado de máquina que utiliza redes neurais com múltiplas camadas (profundas) para extrair e aprender representações complexas de dados.

    \item[\textbf{CNN (Convolutional Neural Network)}] Rede neural especializada na processamento de imagens, utilizada para extrair características visuais através de operações de convolução.

    \item[\textbf{Dataset}] Conjunto de dados utilizado para treinar, validar e testar modelos de aprendizado de máquina.

    \item[\textbf{Deep Learning}] Ver Aprendizado Profundo.

    \item[\textbf{FPS (Frames Per Second)}] Número de quadros de vídeo capturados ou processados por segundo, medida de taxa de amostragem em vídeo.

    \item[\textbf{Gesto Dinâmico}] Movimento realizado com as mãos, braços ou expressão facial que possui significado na Língua Brasileira de Sinais, caracterizado pela mudança de posição ao longo do tempo.

    \item[\textbf{Libras}] Sigla para Língua Brasileira de Sinais, idioma natural utilizado pela comunidade surda brasileira.

    \item[\textbf{Landmark}] Ponto-chave (coordenada tridimensional) que representa uma articulação ou ponto de interesse no corpo humano, extraído por algoritmos de detecção de pose.

    \item[\textbf{Latência}] Intervalo de tempo decorrido entre a execução de um sinal e a exibição da tradução resultante na interface do sistema.

    \item[\textbf{LSTM (Long Short-Term Memory)}] Tipo de rede neural recorrente capaz de aprender dependências de longo prazo em sequências temporais de dados.

    \item[\textbf{MediaPipe}] Framework de código aberto desenvolvido pelo Google para detectar e rastrear landmarks de pose em tempo real a partir de vídeo.

    \item[\textbf{Modelo de Classificação}] Algoritmo treinado que recebe dados de entrada e produz predições de categoria ou classe como saída.

    \item[\textbf{OpenCV}] Biblioteca de código aberto para processamento de imagens e visão computacional, amplamente utilizada em projetos de análise de vídeo.

    \item[\textbf{Pose Estimation}] Tarefa de visão computacional que consiste em localizar e rastrear as articulações do corpo humano em imagens ou vídeos.

    \item[\textbf{RNN (Recurrent Neural Network)}] Rede neural projetada para processar sequências de dados, mantendo memória de informações anteriores.

    \item[\textbf{TensorFlow}] Plataforma de código aberto para aprendizado de máquina desenvolvida pelo Google, utilizada para construir e treinar modelos de redes neurais.

    \item[\textbf{Visão Computacional}] Área da Inteligência Artificial dedicada a capacitar computadores a interpretar e compreender informações visuais a partir de imagens e vídeos.

    \item[\textbf{Vocabulário Acadêmico}] Conjunto de sinais em Libras selecionados para o projeto, relacionados a contextos e atividades comuns em ambientes universitários, como biblioteca, prova e secretaria.

\end{description}